Table~\ref{tab:closedloop} adds onboard sensing that sees every face within the radius, an idealised bound. With $R=20$\,m the rover reaches the goal without touching a face on {{VAL:cl_post_E1_I30_none_20}} of the pairs when its 30\,m prior holds only the slip model, and on {{VAL:cl_post_E1_I30_GBM_20}} with boosting as prior ({{VAL:cl_post_E1_I30_GBM_20_d}} points, interval {{VAL:cl_post_E1_I30_GBM_20_lo}} to {{VAL:cl_post_E1_I30_GBM_20_hi}}); the face map prior gives {{VAL:cl_post_E1_I30_FACEMAP_20}}. No route between E1 endpoints touched a face and no pair used up its replan budget; every failure is a pair left without a route, either before the first move, when blocked cells and faces sensed at the start close every way inside the window (the same pairs for every prior: {{VAL:fb_E1_I30_none_stuckatstart}} of {{VAL:fb_E1_I30_none_pairs}} pairs at 30\,m, {{VAL:fb_E1_I10_none_stuckatstart}} at 10\,m, which is why the 10\,m prior does no better here), or after sensed faces closed the remaining ways. A better prior mainly saves replanning ({{VAL:cl_post_E1_I30_none_20_replans}} replans per pair with the slip model prior, {{VAL:cl_post_E1_I30_FACEMAP_20_replans}} with the face map). With E2 endpoints {{VAL:cl_post_E2_I30_none_20_unsafe}} of the pairs touch a face whatever the prior. A slip face lies behind a crest, out of sight of a rover climbing the windward slope (Section~\ref{sec:posthoc}).
